\documentclass[11pt,oneside,openany]{book}

\usepackage[T1]{fontenc}
\usepackage{lmodern}
\usepackage[letterpaper,margin=1.10in,includeheadfoot]{geometry}
\usepackage{amsmath,amssymb,amsthm,mathtools}
\usepackage{microtype}
\usepackage{xcolor}
\usepackage{titlesec}
\usepackage{fancyhdr}
\usepackage{needspace}
\usepackage{hyperref}

\ifdefined\pdftrailerid
  \pdftrailerid{MathGenome Advanced Finite Group Representation Theory}
\fi

\definecolor{ink}{HTML}{172923}
\definecolor{deepgreen}{HTML}{2B5242}
\definecolor{rulegreen}{HTML}{8A9F92}
\hypersetup{
  colorlinks=true,linkcolor=deepgreen,urlcolor=deepgreen,citecolor=deepgreen,
  pdftitle={Advanced Finite Group Representation Theory},
  pdfsubject={The Symmetric Group and Schur-Weyl Duality},
  pdfauthor={MathGenome},bookmarksnumbered=true
}
\setcounter{tocdepth}{1}
\setcounter{secnumdepth}{2}
\setlength{\headheight}{15pt}
\setlength{\emergencystretch}{2em}
\raggedbottom

\titleformat{\chapter}[display]
  {\normalfont\color{ink}}{\color{deepgreen}\large\scshape Chapter \thechapter}{9pt}
  {\huge\bfseries}
\titlespacing*{\chapter}{0pt}{4pt}{25pt}
\titleformat{\section}{\normalfont\Large\bfseries\color{ink}}{\thesection}{0.7em}{}

\pagestyle{fancy}
\fancyhf{}
\fancyhead[L]{\small\sffamily\color{deepgreen}MathGenome}
\fancyhead[R]{\small\nouppercase{\leftmark}}
\fancyfoot[C]{\small\thepage}
\renewcommand{\headrulewidth}{0.35pt}
\renewcommand{\chaptermark}[1]{\markboth{\thechapter.\ #1}{}}
\fancypagestyle{plain}{%
  \fancyhf{}\fancyfoot[C]{\small\thepage}\renewcommand{\headrulewidth}{0pt}}

\newtheorem{theorem}{Theorem}[chapter]
\newtheorem{lemma}[theorem]{Lemma}
\newtheorem{proposition}[theorem]{Proposition}
\newtheorem{corollary}[theorem]{Corollary}
\theoremstyle{definition}
\newtheorem{definition}[theorem]{Definition}
\newtheorem{example}[theorem]{Example}
\theoremstyle{remark}
\newtheorem{remark}[theorem]{Remark}
\allowdisplaybreaks[1]

\begin{document}
\frontmatter
\hypersetup{pageanchor=false}
\begin{titlepage}
\thispagestyle{empty}
\thispdfpagelabel{Cover}
\begin{flushleft}
\vspace*{0.65in}
{\large\scshape\color{deepgreen}MathGenome\par}
\vspace{0.35in}
{\color{rulegreen}\rule{0.94\textwidth}{0.65pt}\par}
\vspace{0.43in}
{\color{ink}\fontsize{36}{42}\selectfont\bfseries
Advanced Finite Group\\[4pt] Representation Theory\par}
\vspace{0.32in}
{\color{deepgreen}\Large The Symmetric Group and Schur--Weyl Duality\par}
\vfill
{\large Mathematical lecture notes\par}
\vspace{0.10in}
{\normalsize MathGenome v0.21.0 dependency baseline\par}
\vspace*{0.45in}
\end{flushleft}
\end{titlepage}
\hypersetup{pageanchor=true}

\tableofcontents

\chapter*{Scope and conventions}
\addcontentsline{toc}{chapter}{Scope and conventions}

A partition organizes both a conjugacy class of the symmetric group
and an irreducible module. Young tableaux connect these two roles.
They describe the construction and branching of the simple modules,
and give bases and character formulas that can be used in explicit
calculations. These modules occur again in tensor powers, alongside
polynomial representations of the general linear group.

Chapter 1 constructs Specht modules from tabloids, proves that they
are exactly the simple complex representations, and compares the
two orders of the Young symmetrizer. Chapter 2 proves the standard
basis and branching rule, and derives the hook-length formula from
the corner recurrence. Chapter 3 uses the Cauchy kernel to recover
characters and relates them to Schur polynomials, induction, and
border-strip deletion, and gives both Jacobi--Trudi formulas.
Chapter 4 computes both centralizers on
\(V^{\otimes n}\), determines which shapes occur, and classifies
homogeneous polynomial representations. The main references for
these two parts of the story are \cite[Chapters 2--4]{sagan} and
\cite[Sections 4.12--4.22]{etingof}.

Unless a coefficient field is specified, all representations and
vector spaces are finite dimensional over \(\mathbb C\). Rows of
a Young diagram are read left to right, with the longest at the
top. Permutations compose right to left. The slot action on tensor
powers is a left action, with output slot \(i\) coming from
slot \(\sigma^{-1}(i)\). Our Hermitian forms are linear in the first
argument. The empty partition labels the trivial module of \(S_0\),
and the zeroth tensor power is \(\mathbb C\).

The proofs draw on the ordinary character and induction theory in
\cite[Chapters 8--9]{algebra-core} and the weighted path determinant
in \cite[Chapter 3]{finite-combinatorics}. The accompanying source
map lists the precise sources, internal dependencies, and applicable
Lean interfaces for every numbered entry.

\mainmatter
\chapter[The symmetric group and Specht modules]{The Symmetric Group and Specht Modules}

Write \(S_n\) for the permutations of \(\{1,\ldots,n\}\), composed from
right to left. All modules in the main argument are complex and
finite-dimensional. We use the group algebra \(A_n=\mathbb C[S_n]\) and
regard its representations as left modules. The sign of \(g\in S_n\)
is written \(\operatorname{sgn}(g)\). The group \(S_0\) is trivial, and
every construction on zero labels uses the empty diagram.

The finite-group results used below are Schur's lemma, Maschke's
theorem and the count of irreducible complex representations by
conjugacy classes. They are established in \cite[Chapter 8]{algebra-core}.
For the count, the finite-group theorem starts from a complete family
of irreducibles. Its characters are an orthonormal basis of the space
of class functions, whose dimension is the number of conjugacy classes.
We apply this count to determine when a constructed family is complete.

\section{Cycle types and permutation modules}

\begin{theorem}[Conjugacy classes and their sizes]
\label{res:A01}
The cycle type of a permutation records all its cycle lengths, including
one for each fixed point. It gives a bijection from conjugacy classes of
\(S_n\) to partitions of \(n\). If \(\nu\) has \(m_j\) parts of size \(j\),
the centralizer of a permutation of type \(\nu\) has order
\[
  z_\nu=\prod_{j\geq 1}j^{m_j}m_j!,
\]
and the corresponding class has \(n!/z_\nu\) elements. At \(n=0\) these
counts are both \(1\).
\end{theorem}

\begin{proof}
Conjugation relabels each orbit of a permutation and preserves its
cyclic order. If two permutations have the same type, pair their cycles
of equal length and choose a starting label in each paired cycle.
Mapping successive labels in the source cycle to successive labels in
its target defines a bijection that conjugates the permutations.
Conversely, for any partition divide the labels into successive blocks of
its part sizes and cycle the labels of each block. This realizes every
partition, with the sole possibility at \(n=0\) being the identity.

To commute with a given permutation \(g\), a permutation must send each
cycle of \(g\) to a cycle of the same length. For cycles of length \(j\)
there are \(m_j!\) choices of their target order and \(j\) choices of
the image of one chosen starting label on each cycle. Commuting with
\(g\) then determines the images of all successive labels. Each such
choice gives a bijection commuting with \(g\), independently for each
length. The class-size formula now follows from orbit--stabilizer for
conjugation. See \cite[\S2.1]{sagan} and \cite[\S4.14]{etingof}.
\end{proof}

A partition \(\lambda\vdash n\) is a weakly decreasing list of positive
integers summing to \(n\). We may append zeros when comparing lists.
Its Young diagram consists of cells \((i,j)\) with \(1\leq j\leq\lambda_i\).
Rows are numbered from top to bottom and columns from left to right.
The transpose partition \(\lambda^\top\) has \(j\)-th part equal to the
height of column \(j\), and \(\ell(\lambda)\) is the number of nonempty rows.

\begin{proposition}[Tableaux, tabloids, and induction]
\label{res:A02}
A \(\lambda\)-tableau \(t\) is a bijective labeling of the diagram of
\(\lambda\) by \(\{1,\ldots,n\}\). Two tableaux give the same tabloid if
the sets of labels in the corresponding rows agree. Write \(\{t\}\)
for the tabloid of \(t\). The row subgroup \(R_t\) and the column
subgroup \(C_t\) permute labels within the rows and columns of \(t\),
respectively, so \(R_t\cap C_t=\{1\}\). The left permutation module
\(M^\lambda\) on the basis of tabloids is isomorphic to
\(\operatorname{Ind}_{R_t}^{S_n}\mathbf 1\), where \(\mathbf 1\) is the
trivial \(R_t\)-module.
\end{proposition}

\begin{proof}
For a fixed tableau \(t\), any other tableau is \(gt\) for a unique
permutation \(g\), where \(g\) replaces each entry \(a\) by \(g(a)\).
The stabilizer of \(\{t\}\) is exactly \(R_t\) because the tabloid
remembers the set of labels in each row. A permutation preserving both
row sets and column sets fixes the label in their intersection at each
cell. Thus \(R_t\cap C_t\) is trivial. Also \(R_t\) is a product of
symmetric groups on the row sets, the Young subgroup for this choice
of labeling.

With the induction model from \cite[Chapter 9]{algebra-core}, send
\(g\otimes 1\) in \(A_n\otimes_{\mathbb C[R_t]}\mathbb C\) to \(\{gt\}\).
The relation \(gr\otimes 1=g\otimes 1\) holds for every \(r\in R_t\),
and the proposed map respects it. As a right \(\mathbb C[R_t]\)-module,
\(A_n\) is a direct sum of the spans \(g\mathbb C[R_t]\), one for each
left coset \(gR_t\). Tensoring makes each summand a line and sends its
generator to the corresponding tabloid. Hence the map is bijective
and intertwines left multiplication by \(S_n\). See
\cite[\S2.1, Theorem 2.1.12]{sagan}.
\end{proof}

For partitions \(\lambda,\mu\vdash n\), write \(\lambda\unrhd\mu\) if
\(\sum_{i=1}^k\lambda_i\geq\sum_{i=1}^k\mu_i\) for every \(k\).
These sums include appended zeros. In particular, dominance in both
directions forces equal parts.

\begin{lemma}[The row--column test]
\label{res:A03}
If \(\lambda\not\unrhd\mu\), then for any \(\lambda\)-tableau \(t\) and
\(\mu\)-tableau \(u\), some column of \(t\) and some row of \(u\) have at
least two labels in common. If \(\lambda=\mu\) and there is no such
intersection, then \(\{u\}=\{ct\}\) for a unique \(c\in C_t\).
\end{lemma}

\begin{proof}
In the absence of an intersection of size at least two, column \(j\)
of \(t\) has at most \(\min(k,\lambda^\top_j)\) of its labels in the
first \(k\) rows of \(u\). Counting all the labels in those rows gives
\[
  \sum_{i=1}^k\mu_i
  \ \leq\ \sum_{j\geq 1}\min(k,\lambda^\top_j)
  \ =\ \sum_{i=1}^k\lambda_i .
\]
The last equality counts the cells in the first \(k\) rows of the
\(\lambda\)-diagram by columns. This proves the first claim. When the
shapes are equal, the sum is an equality for each \(k\), so every
column attains its individual upper bound. Taking the differences of
these bounds at consecutive \(k\) shows that each column \(j\) supplies
exactly one label to row \(i\) of \(u\) if \(i\leq\lambda^\top_j\),
and none otherwise. Within each column, send the label of \(t\) in
cell \((i,j)\) to that column's unique label in row \(i\) of \(u\).
The resulting \(c\in C_t\) has \(\{ct\}=\{u\}\). Two such choices
could differ only by an element of \(C_t\cap R_t\), which is trivial.
Compare \cite[\S\S2.2--2.3]{sagan}.
\end{proof}

\section{Polytabloids and irreducibility}

\begin{proposition}[Column alternation and the Specht module]
\label{res:A04}
For a tableau \(t\), put
\[
  b_t=\sum_{c\in C_t}\operatorname{sgn}(c)c\in A_n,
  \qquad e_t=b_t\{t\}\in M^\lambda .
\]
The vector \(e_t\) is called a polytabloid. Its coefficient at \(\{t\}\)
is \(1\), and \(g e_t=e_{gt}\) for all \(g\in S_n\).
The Specht module \(S^\lambda\) is the span of the polytabloids of shape
\(\lambda\). It is a nonzero submodule generated by any one \(e_t\).
For \(c\in C_t\), \(e_{ct}=\operatorname{sgn}(c)e_t\), and
\(b_t e_t=|C_t|e_t\).
\end{proposition}

\begin{proof}
By Proposition~\ref{res:A02}, only the identity in \(C_t\) fixes
\(\{t\}\); in fact the \(\{ct\}\) for \(c\in C_t\) are distinct.
Since \(C_{gt}=gC_tg^{-1}\) and sign is invariant under conjugation,
\(b_{gt}=gb_tg^{-1}\), which gives \(ge_t=e_{gt}\). The action on
tableaux is transitive, so the orbit of any \(e_t\) generates the
indicated nonzero module. Multiplying \(b_t\) on either side by
\(c\in C_t\) gives \(\operatorname{sgn}(c)b_t\) by reindexing the sum.
Consequently \(c e_t=\operatorname{sgn}(c)e_t\), and each signed
summand of \(b_t\) acts on \(e_t\) as the identity. This yields
\(b_t e_t=|C_t|e_t\). See \cite[\S2.3]{sagan}.
\end{proof}

Give \(M^\lambda\) the positive definite Hermitian inner product
with orthonormal tabloid basis, linear in the first variable. The
permutation action preserves this inner product.

\begin{theorem}[Submodule theorem]
\label{res:A05}
For a tableau \(t\) of shape \(\lambda\), the action of \(b_t\) on
\(M^\lambda\) is \(b_t x=\langle x,e_t\rangle e_t\).
Every submodule \(U\subseteq M^\lambda\) contains \(S^\lambda\) or
is contained in \((S^\lambda)^\perp\). In particular \(S^\lambda\)
is irreducible.
\end{theorem}

\begin{proof}
First consider the action on an arbitrary tabloid \(\{u\}\), possibly
of a different shape. If some row of \(u\) contains two labels in a
column of \(t\), the transposition of these labels belongs to \(C_t\)
and fixes \(\{u\}\). Pairing terms in \(b_t\) by right multiplication
by this transposition gives \(b_t\{u\}=0\). In the equal-shape case
with no such collision, Lemma~\ref{res:A03} gives
\(\{u\}=c\{t\}\) for \(c\in C_t\). The product rule for \(b_t c\)
then gives \(b_t\{u\}=\operatorname{sgn}(c)e_t\).

Thus the image on \(M^\lambda\) is the line \(\mathbb C e_t\). The
adjoint of each basis permutation is its inverse, and
\(\operatorname{sgn}(c^{-1})=\operatorname{sgn}(c)\), so \(b_t\)
is self-adjoint. If \(b_t x=\alpha e_t\), taking the inner product
against \(\{t\}\) gives
\[
  \alpha=\langle b_t x,\{t\}\rangle
        =\langle x,b_t\{t\}\rangle=\langle x,e_t\rangle .
\]
Here the first equality uses the coefficient \(1\) from
Proposition~\ref{res:A04}. If some vector \(x\in U\) pairs nontrivially
with any \(e_t\), applying \(b_t\) puts that \(e_t\) in \(U\), and the
orbit of \(e_t\) gives \(S^\lambda\subseteq U\). Otherwise \(U\)
is perpendicular to every \(e_t\) and hence to \(S^\lambda\).
Apply this dichotomy to a submodule of \(S^\lambda\). In the second
case it is zero by positive definiteness, while in the first it is all
of \(S^\lambda\). This is the argument of \cite[\S2.4]{sagan}.
\end{proof}

\begin{proposition}[Dominance and diagonal multiplicity]
\label{res:A06}
If \(\operatorname{Hom}_{S_n}(S^\lambda,M^\mu)\neq 0\), then
\(\lambda\unrhd\mu\). The space
\(\operatorname{Hom}_{S_n}(S^\lambda,M^\lambda)\) has dimension \(1\).
\end{proposition}

\begin{proof}
An intertwiner \(f:S^\lambda\to M^\mu\) is determined by \(f(e_t)\)
for a fixed \(\lambda\)-tableau \(t\), since \(e_t\) generates.
Moreover
\[
  b_t f(e_t)=f(b_t e_t)=|C_t| f(e_t),
\]
so \(f(e_t)\) belongs to the image of \(b_t\). If
\(\lambda\not\unrhd\mu\), the first part of Lemma~\ref{res:A03}
and the transposition pairing in Theorem~\ref{res:A05} show
that \(b_t\) kills every basis tabloid of \(M^\mu\).
If \(\lambda=\mu\), Theorem~\ref{res:A05} gives the one-dimensional
image \(\mathbb C e_t\), and the inclusion \(S^\lambda\hookrightarrow
M^\lambda\) gives a nonzero intertwiner. See
\cite[\S2.4, Proposition 2.4.5]{sagan}.
\end{proof}

\begin{theorem}[Classification over \(\mathbb C\)]
\label{res:A07}
The modules \(S^\lambda\), for \(\lambda\vdash n\), form a complete
family of pairwise nonisomorphic irreducible \(S_n\)-modules.
For \(n\geq 1\), \(S^{(n)}\) is trivial and \(S^{(1^n)}\) is the
sign module. At \(n=0\), \(S^\varnothing\) is the trivial line.
\end{theorem}

\begin{proof}
Theorem~\ref{res:A05} proves irreducibility. If
\(S^\lambda\simeq S^\mu\), then composing an isomorphism and its
inverse with the respective inclusions in \(M^\mu\) and \(M^\lambda\)
gives nonzero intertwiners. Proposition~\ref{res:A06} yields both
dominance inequalities, so \(\lambda=\mu\).
Theorem~\ref{res:A01} gives a conjugacy class for each partition,
so the number constructed is already the number of isomorphism
classes given by the finite-group count.

For one row there is one tabloid and every permutation fixes it.
For one column, \(C_t=S_n\), and Proposition~\ref{res:A04} makes
the orbit of \(e_t\) the sign line. The empty diagram has one tabloid
and the only permutation fixes it. The classification is
\cite[\S2.4, Theorem 2.4.6]{sagan}; compare
\cite[\S4.12, Theorem 4.36]{etingof}.
\end{proof}

\section{Young operators in the group algebra}

\begin{proposition}[Left ideals and normalized Young idempotents]
\label{res:A08}
Fix a tableau \(t\) of shape \(\lambda\) and put
\[
  a_t=\sum_{r\in R_t}r,\qquad c_t=b_t a_t,\qquad c'_t=a_t b_t .
\]
The map \(\{gt\}\mapsto ga_t\) identifies \(M^\lambda\) with
\(A_n a_t\), and its restriction identifies \(S^\lambda\) with
\(A_n c_t\). With \(q_\lambda=n!/\dim S^\lambda\),
\[
  c_t^2=q_\lambda c_t,\qquad (c'_t)^2=q_\lambda c'_t .
\]
Thus \(c_t/q_\lambda\) and \(c'_t/q_\lambda\) are idempotents, and
the left ideals \(A_n c_t\) and \(A_n c'_t\) are isomorphic.
\end{proposition}

\begin{proof}
The map on tabloids is well-defined since \(ra_t=a_t\) for \(r\in R_t\).
The sum \(ga_t\) has support on \(gR_t\), and sums on distinct cosets
have disjoint supports. Hence the map identifies the tabloid basis
with a basis of \(A_n a_t\) and is equivariant. It sends \(e_t\)
to \(b_t a_t=c_t\), which generates the image \(A_n c_t\).

Right multiplication by \(c_t\) on \(A_n\) is a left-module map with
image \(A_n c_t\). Its restriction to that image is multiplication
by some scalar \(q\) by Schur's lemma, since \(A_n c_t\simeq S^\lambda\)
is irreducible. In the group basis its trace is \(n!\) times the
coefficient of \(1\) in \(c_t\). That coefficient is \(1\): an equality
\(cr=1\) for \(c\in C_t\), \(r\in R_t\) implies \(c=r=1\).
The map induces zero on the quotient \(A_n/(A_n c_t)\), so its trace
also equals \(q\dim S^\lambda\). Thus \(q=q_\lambda\), and applying
the restriction to \(c_t\) gives \(c_t^2=q_\lambda c_t\).
The linear map \(g\mapsto g^{-1}\) reverses products and fixes \(a_t\)
and \(b_t\), so it takes the square relation for \(c_t\) to the
one for \(c'_t\).

Right multiplication by \(b_t\) maps \(A_n c_t\) into \(A_n c'_t\)
because \(c_t b_t=b_t c'_t\). Right multiplication by \(a_t\) maps
back because \(c'_t a_t=a_t c_t\). The compositions are right
multiplication by \(c_t\) and \(c'_t\) on the respective ideals,
each equal to scalar \(q_\lambda\) by the square relations. These
maps are therefore inverse up to a nonzero scalar. Our column-sum
then row-sum order is the \(b_t a_t\) model compared with \(a_t b_t\)
in \cite[Lemma 4.42 and Problem 4.68(a)]{etingof}.
\end{proof}

\begin{corollary}[Transpose and sign twist]
\label{res:A09}
For every \(\lambda\vdash n\),
\(S^{\lambda^\top}\simeq S^\lambda\otimes\operatorname{sgn}\).
The regular module of \(S_2\) over \(\mathbb F_2\) has a fixed line
with no invariant complement. It illustrates the role of invertibility
of the group order in Maschke's theorem.
\end{corollary}

\begin{proof}
The algebra automorphism \(\theta\) defined by
\(\theta(g)=\operatorname{sgn}(g)g\) is an involution. For a left
ideal \(I\), the linear map \(x\mapsto\theta(x)\) intertwines the
sign-twisted action on \(I\) with the ordinary action on \(\theta(I)\),
since \(\theta(\operatorname{sgn}(g)gx)=g\theta(x)\).
The transposed tableau \(t^\top\) interchanges the row and column
subgroups. Hence \(\theta(b_t)=a_{t^\top}\),
\(\theta(a_t)=b_{t^\top}\), and \(\theta(c_t)=c'_{t^\top}\).
Proposition~\ref{res:A08} identifies \(A_n c'_{t^\top}\) with
\(S^{\lambda^\top}\), proving the sign twist formula.

For \(S_2=\{1,s\}\) over \(\mathbb F_2\), the regular module has
basis \(1,s\), and \(s\) exchanges the two basis vectors. The fixed
space is the line spanned by \(1+s\). An invariant complement would
have dimension \(1\), and on it \(s\) would act by a nonzero scalar
in \(\mathbb F_2\), necessarily \(1\). Such a complement would
therefore lie in the fixed line, which is impossible.
Compare \cite[\S3.1 and Problem 4.68(b)]{etingof}.
\end{proof}

\chapter{Standard Tableaux and Branching}

The spanning set of all polytabloids from Result~\ref{res:A04} is
usually much larger than the dimension of the Specht module. The
standard polytabloids give a basis whose labels track restriction to
\(S_{n-1}\). The same labels yield an integer recurrence for the
dimension, which we solve by a determinant and then by hook lengths.
Our conventions for diagrams and tabloids are those of
Results~\ref{res:A02} and~\ref{res:A04}.

\section{The standard polytabloid basis}

A tableau \(t\) of shape \(\lambda\vdash n\) is \emph{standard} if its
labels \(1,\ldots,n\) increase strictly to the right along each row and
down each column. Call it \emph{column-increasing} if it satisfies only
the column condition. The empty diagram has a single tableau, which is
standard.

\begin{proposition}[Independence of standard polytabloids]\label{res:B01}
The vectors \(e_t\) for standard \(t\) are linearly independent in
\(M^\lambda\).
\end{proposition}

\begin{proof}
For a tabloid \(u\), let \(N_{k,m}(u)\) count the labels at most \(m\)
in its first \(k\) rows. These counts for all \(k,m\) determine \(u\):
the difference \(N_{k,m}(u)-N_{k,m-1}(u)\) tells whether the label
\(m\) is among the first \(k\) rows. Thus
\[
  u\preceq v\quad\Longleftrightarrow\quad
  N_{k,m}(u)\le N_{k,m}(v)\quad\text{for every } k,m
\]
is a partial order on tabloids of the given shape.

For a column-increasing \(t\), the entries of each column retained
in the first \(k\) rows are the smallest entries of that column in
the required number. Reordering any column cannot increase the count
of entries at most \(m\) in these positions. Hence
\(\{ct\}\preceq\{t\}\) for each \(c\in C_t\). The coefficient of
\(\{t\}\) in \(e_t\) is \(1\) by Result~\ref{res:A04}, and all other
terms are strictly smaller in the order just defined. Distinct
standard tableaux have distinct tabloids, since their row order can
be recovered by sorting each row.

If a finite linear combination of standard polytabloids had sum zero
and some coefficient were nonzero, take a maximal tabloid among the
leading tabloids with nonzero coefficient. No other summand contributes
at it, by the order property and distinctness. Its coefficient in the
sum is therefore nonzero, a contradiction. This is the triangular
argument of \cite[Section 2.5, Proposition 2.5.9]{sagan}.
\end{proof}

\begin{proposition}[Garnir relations and straightening]\label{res:B02}
Suppose \(t\) is column-increasing and has a row descent at adjacent
columns, with entries \(a_i>b_i\) in row \(i\). Write the two columns
as \(a_1<\cdots<a_p\) and \(b_1<\cdots<b_q\), and set
\[
  A=\{a_i,\ldots,a_p\},\qquad B=\{b_1,\ldots,b_i\}.
\]
For a set of labels \(U\), write \(S_U\) for the group that permutes
\(U\) and fixes all other labels. If \(P\) represents the left cosets
of \(S_A\times S_B\) in \(S_{A\cup B}\) and contains the identity, then
\[
  \sum_{\pi\in P}\operatorname{sgn}(\pi)e_{\pi t}=0.
\]
Every polytabloid is a linear combination of standard polytabloids.
\end{proposition}

\begin{proof}
Write \(U=A\cup B\) and \(H=S_A\times S_B\). Every term \(\{ct\}\)
of \(e_t\), with \(c\in C_t\), has the labels of \(U\) within the
same two columns. Those columns occupy at most \(p\) rows, since
\(q\le p\), while \(\lvert U\rvert=p+1\). Some two labels of \(U\)
therefore lie in the same row of \(ct\). Their transposition \(\tau\)
fixes \(\{ct\}\), and the terms indexed by \(\sigma\) and
\(\sigma\tau\) in the signed sum over \(S_U\) cancel on this tabloid.
The signed sum over \(S_U\) consequently annihilates \(e_t\).

The subgroup \(H\) is contained in \(C_t\), so \(h e_t\)
equals \(\operatorname{sgn}(h)e_t\) by Result~\ref{res:A04}. Factor the
signed sum over \(S_U\) using its disjoint cosets \(\pi H\). The inner
signed sum over \(H\) acts by \(\lvert H\rvert\), and dividing the
result by this nonzero integer proves the relation.

To straighten, first sort the columns of any tableau. A column
permutation acts on its polytabloid by its sign, so this step preserves
the span. Order the column-increasing tableaux by lexicographic order
on the word obtained by reading each column from top to bottom and
the columns from left to right. A column-increasing tableau that is
not standard has a descent as in the statement. If \(\pi\) represents
a nonidentity coset, then \(\pi A\ne A\). At least one label of \(A\)
is replaced in the left column by a label of \(B\). All labels of
\(B\) are smaller than every label of \(A\), because \(b_i<a_i\).
All other entries of this column are unchanged, and the earlier
columns are unchanged. After sorting, the first difference in the
left column is a smaller label, since the smallest label in the
symmetric difference of the two column sets belongs to the new set.
The column word has strictly decreased.

Isolate \(e_t\) in the relation and sort every other term, keeping the
sign of each sorting permutation. All resulting polytabloids have
smaller words, so induction over the finite order expresses each one
in the span of the standard polytabloids. The relation and
straightening are the column form of the argument in
\cite[Section 2.6, Proposition 2.6.3 and Theorem 2.6.4]{sagan}.
\end{proof}

\begin{theorem}[Standard basis and rational model]\label{res:B03}
The standard polytabloids form a basis of \(S^\lambda\). In
particular, \(\dim_{\mathbb C} S^\lambda=f^\lambda\), where
\(f^\lambda\) is the number of standard tableaux of shape \(\lambda\).
The span of all polytabloids in the rational tabloid module is a
rational \(S_n\)-module whose complexification is \(S^\lambda\).
\end{theorem}

\begin{proof}
Every generator of \(S^\lambda\) straightens by
Proposition~\ref{res:B02}, and Proposition~\ref{res:B01} gives
independence. The tabloid basis and every polytabloid have integral
coordinates. The signed coset relation and its use in straightening
hold over \(\mathbb Q\), and the same triangular proof gives
independence over \(\mathbb Q\). The resulting rational and complex
bases are indexed by the same standard tableaux. The action on
tabloids is given by permutations, so the scalar-extension map on
these bases is an isomorphism of \(S_n\)-modules
\cite[Theorem 2.6.5]{sagan}, \cite[Corollary 4.38]{etingof}.
\end{proof}

\section{Removing and adding one box}

A removable box of a nonempty diagram is a cell whose deletion
leaves a diagram. These are the last cells of rows \(r\) with
\(\lambda_r>\lambda_{r+1}\), where trailing parts are zero. For each
\(n\ge 1\), view \(S_{n-1}\) as the subgroup of \(S_n\) fixing the
label \(n\).

\begin{proposition}[Restriction filtration]\label{res:B04}
For \(n\ge 1\), let \(c_1,\ldots,c_s\) be the removable boxes of
\(\lambda\vdash n\), listed by increasing row numbers
\(r_1,\ldots,r_s\). Put
\(\mu_i=\lambda\setminus c_i\). For \(1\le i\le s\) let \(F_i\)
be the span of standard \(e_t\) with \(n\) in one of
\(c_1,\ldots,c_i\) and put \(F_0=0\). Then \(F_i\) are
\(S_{n-1}\)-submodules, and deleting \(n\) from the \(i\)-th row
choice induces \(F_i/F_{i-1}\simeq S^{\mu_i}\). Thus the full
restriction has a filtration with one successive quotient for
each removable box.
\end{proposition}

\begin{proof}
For each \(j\), define \(D_j:M^\lambda\to M^{\mu_j}\) on tabloids.
If \(n\) is in row \(r_j\), delete it from that row. Otherwise send
the tabloid to zero. This depends only on row sets and commutes with
all permutations fixing \(n\). The target row lengths are those of
\(\mu_j\) because \(c_j\) was removable. If this deletion empties the
last row, omit that row.

In a standard tableau \(t\), \(n\) has no neighbor to its right or
below, since such a neighbor would have a greater label. Suppose \(n\)
is in \(c_a\). In every term of \(e_t\), a column permutation can only
move \(n\) upwards within its column, since \(c_a\) is the bottom cell
of the column. Thus \(D_j(e_t)=0\) for \(j>a\). For \(j=a\), only
the column permutations fixing \(n\) contribute. Restricting them
to the other labels identifies them with the column group of
\(t\setminus n\) and preserves their signs. Therefore
\[
  D_a(e_t)=e_{t\setminus n}.
\]
Deleting \(n\) from standard tableaux with \(n\) in \(c_a\) gives
a bijection onto the standard tableaux of shape \(\mu_a\), with
inverse given by placing \(n\) in \(c_a\).

For each \(i\) we now have the equality of vector subspaces
\[
  F_i=S^\lambda\cap\bigcap_{j>i}\ker D_j .
\]
One inclusion follows from the vanishing just proved. If a vector on
the right had a nonzero coefficient at a corner with index greater
than \(i\), take the largest such index \(a\) in its standard-basis
expansion. Applying \(D_a\) gives a nonzero combination of distinct
standard polytabloids of \(S^{\mu_a}\); terms from all smaller corners
vanish, and no term from a larger corner occurs. This contradicts
membership in the kernel. The same argument includes \(i=0\).
The characterization by kernels proves invariance under \(S_{n-1}\).
On \(F_i\), the map \(D_i\) has image \(S^{\mu_i}\) and kernel
\(F_{i-1}\), by its effect on the standard bases. This proves the
quotient assertion \cite[Section 2.8]{sagan}.
\end{proof}

\begin{theorem}[Branching rule]\label{res:B05}
For \(\lambda\vdash n\) with \(n\ge 1\) and \(\nu\vdash m\)
with \(m\ge 0\), one has
\[
  \begin{aligned}
  \operatorname{Res}^{S_n}_{S_{n-1}} S^\lambda
    &\simeq \bigoplus_{\mu\text{ obtained by removing one box}} S^\mu,\\
  \operatorname{Ind}^{S_{m+1}}_{S_m} S^\nu
    &\simeq \bigoplus_{\eta\text{ obtained by adding one box}} S^\eta .
  \end{aligned}
\]
The diagrams in the two sums are obtained from \(\lambda\) and
\(\nu\), respectively. Each summand occurs once.
\end{theorem}

\begin{proof}
Maschke's theorem over \(\mathbb C\) splits the restriction filtration
of Proposition~\ref{res:B04}. Two removable boxes in distinct rows
leave distinct sequences of row lengths and therefore distinct
irreducibles by Result~\ref{res:A07}. For the induction statement
use Frobenius reciprocity and Schur's lemma, in the finite-group
forms of \cite[Chapters 8--9]{algebra-core}. For each
\(\eta\vdash m+1\), the multiplicity of \(S^\eta\) in the induced
module is
\[
  \dim_{\mathbb C}\operatorname{Hom}_{S_m}
      \bigl(S^\nu,\operatorname{Res}^{S_{m+1}}_{S_m}S^\eta\bigr).
\]
The restriction formula gives dimension \(1\) precisely when
\(\nu\) results from deleting a corner of \(\eta\). Otherwise the
dimension is zero. All irreducibles are covered by
Result~\ref{res:A07}, so complete reducibility gives the stated sum
\cite[Theorem 2.8.3]{sagan}. For \(n=1\) the restriction deletes the
only box and gives the trivial \(S_0\)-module. Induction from \(S_0\)
to \(S_1\) adds that box and gives the trivial \(S_1\)-module.
\end{proof}

\begin{corollary}[Corner recurrence]\label{res:B06}
The numbers of standard tableaux satisfy
\[
   f^{\varnothing}=1,\qquad
   f^\lambda=\sum_{c\text{ removable in }\lambda}f^{\lambda\setminus c}
   \quad(\lambda\ne\varnothing).
\]
\end{corollary}

\begin{proof}
Group standard tableaux according to the position of the greatest
label. As shown in the proof of Proposition~\ref{res:B04}, that
position must be a corner, and deleting the label and filling it
back are inverse bijections. Alternatively, the recurrence for a
nonempty diagram follows by taking dimensions in the restriction
formula and using Theorem~\ref{res:B03}
\cite[Lemma 2.8.2 and Theorem 2.8.3]{sagan}.
\end{proof}

\section{A determinant and the hook product}

\begin{theorem}[Frobenius--Young dimension determinant]\label{res:B07}
Fix \(r\ge \ell(\lambda)\) and pad \(\lambda\vdash n\) with zeros
through row \(r\). For integer \(a\) put \(q(a)=1/a!\) if
\(a\ge 0\) and \(q(a)=0\) otherwise. Then
\[
  f^\lambda=n!\det\bigl[q(\lambda_i-i+j)\bigr]_{1\le i,j\le r}.
\]
The determinant of the empty matrix is \(1\).
\end{theorem}

\begin{proof}
Let \(D(\lambda)\) be the displayed determinant without the factor
\(n!\). For \(\lambda=\varnothing\) the matrix is upper triangular
with all diagonal entries \(1\). For \(\lambda\ne\varnothing\),
form \(D_i\) by reducing \(\lambda_i\) by \(1\) in the determinant
formula, with no requirement yet that the result is a partition.
If \(i<r\) and \(\lambda_i=\lambda_{i+1}\), the new \(i\)-th row
equals row \(i+1\). If \(i=r\) and \(\lambda_r=0\), the new last row
is zero. These are exactly the invalid deletions, and they give
\(D_i=0\).

For every integer \(a\), \(q(a-1)=a q(a)\), including \(a=0\) and
negative \(a\) by the zero convention. In the determinant expansion,
the signed product belonging to a permutation \(\sigma\) is
therefore multiplied in \(D_i\) by \(\lambda_i-i+\sigma(i)\).
Summing over \(i\) multiplies that product by \(n\), since
\(\sum_i\sigma(i)=\sum_i i\). It follows that
\[
    n D(\lambda)=\sum_i D_i
      =\sum_{c\text{ removable in }\lambda}D(\lambda\setminus c).
\]
Thus \(n!D(\lambda)\) and \(f^\lambda\) have the same initial value
and satisfy the recurrence of Corollary~\ref{res:B06}. Induction on \(n\)
proves the formula for every fixed \(r\) and every diagram of at
most \(r\) rows. This is a recurrence proof of the determinant
in \cite[Theorem 3.11.1]{sagan}.
\end{proof}

For a box \(c=(i,j)\) of \(\lambda\), its \emph{hook} consists of the
box itself, the boxes to its right in its row, and the boxes below it
in its column. If \(\lambda'_j\) is the height of column \(j\), its
length is
\[
    h(i,j)=\lambda_i-j+\lambda'_j-i+1.
\]

\begin{theorem}[Hook-length formula]\label{res:B08}
For every \(\lambda\vdash n\),
\[
      \dim_{\mathbb C} S^\lambda
        = f^\lambda
        = \frac{n!}{\displaystyle\prod_{c\in\lambda}h(c)} .
\]
The empty product is \(1\). In particular the diagram with one row
and the diagram with one column each give dimension one.
\end{theorem}

\begin{proof}
Keep the number \(r\) of padded rows fixed and put
\(b_i=\lambda_i+r-i\). These are strictly decreasing nonnegative
integers. Multiply row \(i\) of the determinant in
Theorem~\ref{res:B07} by \(b_i!\). The \(j\)-th entry becomes
the falling factorial
\[
    (b_i)_{r-j}=b_i(b_i-1)\cdots(b_i-r+j+1),
\]
with \((b_i)_0=1\). If \(b_i<r-j\), the product vanishes, as does
the original entry by the zero convention. Each column consists of
values of a monic polynomial of degree \(r-j\). Starting with the
columns of lowest degree, add combinations of lower degree columns
to replace each one by the corresponding power column. The result
is the Vandermonde determinant with columns of degrees
\(r-1,r-2,\ldots,0\), whose value is
\(\prod_{i<k}(b_i-b_k)\). For example, its vanishing at coincident
variables supplies all these factors, and their degree equals
the degree of the determinant. The coefficient of
\(b_1^{r-1}b_2^{r-2}\cdots b_r^0\) is \(1\) in both expressions.
Consequently
\[
  D(\lambda)=\frac{\prod_{i<k}(b_i-b_k)}{\prod_i b_i!}.
\]

We next identify the hook lengths in row \(i\). They are distinct
positive integers at most \(b_i\): moving one box right strictly
decreases \(\lambda_i-j\) and cannot increase \(\lambda'_j\), and
the greatest possible value at \(j=1\) is \(\lambda_i+r-i\).
For each \(k>i\) the integer \(b_i-b_k\) lies in the same interval,
and these excluded integers are distinct. None is a hook length
in that row. Indeed,
\[
   h(i,j)-(b_i-b_k)=\lambda_k-j+\lambda'_j-k+1.
\]
If \(\lambda_k\ge j\), then \(\lambda'_j\ge k\), so this difference
is positive. Otherwise \(\lambda'_j<k\), so it is negative. The
\(\lambda_i\) hook lengths and the \(r-i\) excluded integers are
therefore disjoint and together number \(b_i\), exhausting
\(\{1,\ldots,b_i\}\). The argument also covers a padded zero row.
Multiplying the row identities gives
\[
   \prod_{c\in\lambda}h(c)
      = \frac{\prod_i b_i!}{\prod_{i<k}(b_i-b_k)}.
\]
The determinant formula and Theorem~\ref{res:B03} now give the
claimed dimension \cite[Theorem 3.10.2]{sagan},
\cite[Theorem 4.53]{etingof}. In a single row or column the hook
lengths are \(n,n-1,\ldots,1\), as required.
\end{proof}

\begin{example}[Small diagrams and the standard representation]\label{res:B09}
For \(n\ge 2\), the sum-zero subspace of the natural point
representation of \(S_n\) is \(S^{(n-1,1)}\). In \(S_3\) this is the
two-dimensional standard representation. The \(S_4\)-modules
\(S^{(2,2)}\) and \(S^{(3,1)}\) have dimensions \(2\) and \(3\).
Their restrictions to \(S_3\) are \(S^{(2,1)}\) and
\(S^{(3)}\oplus S^{(2,1)}\), respectively.
\end{example}

\begin{proof}
For \((n-1,1)\), the bottom-row label uniquely determines a tabloid.
Thus \(M^{(n-1,1)}\) is the point representation
\(\mathbb C^n\), with basis vectors \(v_1,\ldots,v_n\).
If the bottom label is \(b\) and the top label in its column is
\(a\), the polytabloid is \(v_b-v_a\): the only nontrivial
permutation of that column interchanges \(a\) and \(b\). Every
distinct pair can occur. These differences span the sum-zero
subspace, since any vector with coordinates \(x_1,\ldots,x_n\)
summing to zero equals \(\sum_{i<n}x_i(v_i-v_n)\).
For \(n=3\), the \((2,1)\) hook lengths are \(3,1,1\), in
agreement with its dimension \(2\).

For \((2,2)\) the hooks by rows are \(3,2\) and \(2,1\).
For \((3,1)\) they are \(4,2,1\) and \(1\). The hook formula gives
\(4!/12=2\) and \(4!/8=3\). The square diagram has only the lower
right corner, with deletion shape \((2,1)\). In \((3,1)\), deleting
the top right box gives \((2,1)\) and deleting the lower box gives
\((3)\). The decompositions follow by Theorem~\ref{res:B05}.
In the opposite direction, adjoining a box to \((2,1)\) yields
\((3,1),(2,2),(2,1,1)\). The dimensions in the induced module sum
to \(3+2+3=8\), where the last \(3\) also follows by transposition
and the sign twist of Result~\ref{res:A09}. This agrees with the
index \(4\) times the original dimension \(2\).
\end{proof}

\chapter[Characters and Schur polynomials]{Characters and Schur Polynomials}

Cycle lengths describe the fixed tabloids of a permutation. An
alternating coefficient extraction then separates the irreducible
characters. The same coefficients relate these characters to symmetric
polynomials and turn induction into multiplication.

\section{Cycle counting and the Frobenius formula}

For \(\nu\vdash n\), write \(m_j(\nu)\) for the multiplicity of \(j\).
Use \(z_\nu=\prod_j j^{m_j(\nu)}m_j(\nu)!\), as in
Result~\ref{res:A01}. In \(r\) variables set
\[
 p_j(x)=\sum_{i=1}^r x_i^j,\qquad
 p_\nu(x)=\prod_{j\ge1}p_j(x)^{m_j(\nu)}.
\]
All products over an empty set are \(1\). The variables in coefficient
extractions commute, and an exponent with a negative component has
coefficient zero.

\begin{proposition}[The Young permutation character]\label{res:C01}
Let \(\mu=(\mu_1,\ldots,\mu_r)\) be a list of nonnegative integers summing
to \(n\). Let \(\xi^\mu\) be the character of the \(S_n\)-module with
basis the ordered decompositions of the labels into rows of these sizes.
For cycle type \(\nu\), one has \(\xi^\mu(\nu)=[x^\mu]p_\nu(x)\).
\end{proposition}

\begin{proof}
This is the tabloid module of Result~\ref{res:A02}, with ordered row
sizes. Reordering the sizes gives an equivariant bijection by
reordering the rows. The trace on the basis counts fixed decompositions.
A row is preserved by the permutation precisely when it is a union of
cycles. A cycle of length \(j\) can therefore be assigned to any row
\(i\), contributing \(x_i^j\). The product records independent assignments
of the individual cycles, and coefficient extraction imposes the sizes
of all rows. See \cite[Section 4.14]{etingof}.
\end{proof}

For \(r\ge1\), let \(\delta=(r-1,r-2,\ldots,0)\). For a vector \(b\)
of \(r\) nonnegative integers define \(a_b(x)=\det[x_i^{b_j}]_{i,j=1}^r\).
Thus \(\Delta(x)=a_\delta(x)=\prod_{i<j}(x_i-x_j)\). The row convention
in this determinant fixes all the signs in the coefficient formula.

\begin{lemma}[Cauchy kernel and alternating coefficients]\label{res:C02}
For \(r\) variables in each alphabet,
\begin{align*}
 \prod_{i,j}(1-x_i y_j)^{-1}
   &=\sum_\nu p_\nu(x)p_\nu(y)/z_\nu,\\
 \Delta(x)\Delta(y)\prod_{i,j}(1-x_i y_j)^{-1}
   &=\sum_{\ell(\nu)\le r}a_{\nu+\delta}(x)a_{\nu+\delta}(y).
\end{align*}
Consequently, the coefficient of \(x^{\lambda+\delta}y^{\mu+\delta}\)
in the second line is \(1\) if \(\lambda=\mu\) and \(0\) otherwise.
\end{lemma}

\begin{proof}
The first identity follows by expanding
\(\exp(\sum_{j\ge1}p_j(x)p_j(y)/j)\): using the \(j\)-th summand \(m\)
times contributes \(1/(j^m m!)\). All operations take place degree by
degree in formal power series.

For the second identity, clear the product of denominators in the
determinant of the matrix \([(1-x_i y_j)^{-1}]\).
The resulting polynomial alternates separately in \(x\) and in \(y\),
and its degree in each individual \(x_i\) or \(y_j\) is at most \(r-1\).
It is thus a constant times \(\Delta(x)\Delta(y)\). To determine the
constant, set \(x_r=y_r=0\) and subtract the last row from each preceding
row in the determinant before clearing denominators. The last column is
now zero except for its last entry \(1\), and the remaining entries are
\(x_i y_j/(1-x_i y_j)\). Factoring out these row and column factors
reduces the constant to the one for \(r-1\), with value \(1\) for \(r=1\).

Expand each entry as \(\sum_{k\ge0}x_i^k y_j^k\). For a finite truncation,
the rectangular factors have entries \(x_i^k\) and \(y_j^k\) in the
complex polynomial ring. Apply the rectangular Cauchy--Binet identity
over its field of fractions. Both sides are polynomials, so the
identity holds in the polynomial ring by injectivity of the embedding.
This gives a sum of products of determinants indexed by distinct
exponents \(b_1>\cdots>b_r\ge0\).
Using decreasing order in both determinants preserves their product.
In any fixed degree the sum stabilizes. Such exponent lists are
exactly \(\nu+\delta\) for the stated partitions. A monomial with
strictly decreasing exponents can occur only in the corresponding
determinant and then has coefficient \(1\). This proves the extraction
assertion. Compare \cite[Section 4.15]{etingof}.
\end{proof}

\begin{theorem}[Frobenius character formula]\label{res:C03}
For \(\lambda,\nu\vdash n\), put \(\chi^\lambda=\chi_{S^\lambda}\).
If \(r\ge\ell(\lambda)\) and \(r\ge1\), then
\[
 \chi^\lambda(\nu)
   =[x^{\lambda+\delta}]\Delta(x)p_\nu(x),
\]
where \(\lambda\) is padded with zeros to length \(r\). For the empty
partition the formula also has value \(1\) with zero variables.
\end{theorem}

\begin{proof}
First choose \(r\ge \max(1,n)\). As a function of \(\nu\), the indicated
coefficient is \(\xi^\lambda\) plus an integer combination of
\(\xi^\mu\) for shapes \(\mu\) strictly dominating \(\lambda\).
Indeed the determinant expansion takes signed coefficients with
exponents \(\lambda+\delta-w\delta\). When these exponents are
nonnegative, sort them in decreasing order to obtain \(\mu\).
The sum of the largest \(k\) exponents is at least the sum of the
first \(k\) exponents in the unsorted list. The latter is at least
\(\sum_{i=1}^k\lambda_i\), since the first \(k\) entries of \(\delta\)
have maximal sum among its \(k\)-element subsets. If \(w\ne1\), some
initial subset changes, so at least one of the inequalities is strict.
Result~\ref{res:C01} identifies the corresponding coefficient with
\(\xi^\mu\).

By Results~\ref{res:A06} and \ref{res:A07}, \(\xi^\lambda\) contains
\(\chi^\lambda\) once and only strictly dominating shapes can also
occur. Thus our candidate list is unitriangular in the irreducible
characters, in any order with dominating shapes first. Its entries are
orthonormal: the inner product of two candidates is extraction at
\((\lambda+\delta,\mu+\delta)\) in the first line of
Result~\ref{res:C02} multiplied by \(\Delta(x)\Delta(y)\).
Here the candidates are real, and \(1/z_\nu\) is the weight of the
class in the character inner product. The second line gives the desired
orthonormality. Reading the unitriangular list from most dominant to
least, orthogonality to all earlier entries forces all off-diagonal
coefficients to vanish. We used the finite-group character
orthogonality established in \cite[Chapter 8]{algebra-core}.

To remove a final zero part, specialize the final variable to zero.
The power sums are unchanged, and
\(\Delta_{r+1}(x,0)=(x_1\cdots x_r)\Delta_r(x)\).
The desired coefficient in \(r+1\) variables has last exponent \(0\)
and all other exponents raised by \(1\); it therefore reduces to the
stated \(r\)-variable coefficient. See \cite[Theorem 4.47]{etingof}.
\end{proof}

\section{Schur polynomials and the characteristic map}

For \(r\) variables let \(h_k\) be the sum of the weights of weakly
increasing lists of length \(k\) in \(\{1,\ldots,r\}\), where a letter
\(i\) has weight \(x_i\). Set \(h_0=1\) and \(h_k=0\) for \(k<0\).
Write \(e_k\) for the corresponding sum over strictly increasing lists,
and put \(e_0=1\).

\begin{theorem}[Three formulas for the Schur polynomial]\label{res:C04}
For a partition \(\lambda\) with at most \(r\) parts,
\[
 s_\lambda(x)
 =\frac{a_{\lambda+\delta}(x)}{\Delta(x)}
 =\det[h_{\lambda_i-i+j}(x)]_{i,j=1}^r
 =\sum_T \prod_{u\in\lambda}x_{T(u)} .
\]
Here \(T\) ranges over the tableaux with entries in
\(\{1,\ldots,r\}\), weakly increasing along rows and strictly
increasing down columns. For \(\ell(\lambda)>r\) the tableau sum and
the Jacobi--Trudi determinant using at least \(\ell(\lambda)\) rows
are zero. The empty partition gives \(1\).
\end{theorem}

\begin{proof}
The quotient is a symmetric polynomial because the alternating
numerator is divisible by every \(x_i-x_j\). Divide the second identity
of Result~\ref{res:C02} by \(\Delta(x)\). Extract the coefficient of
\(y^{\lambda+\delta}\) in the resulting equality. On the sum side this
gives \(s_\lambda(x)\). On the other side, expand \(\Delta(y)\) and use
\(\prod_{i}(1-x_i y_j)^{-1}=\sum_{k\ge0}h_k(x)y_j^k\) for each \(j\).
The coefficient of the \(i\)-th variable is \(h_{\lambda_i-i+j}\) in
the term using the \(j\)-th exponent of \(\delta\); the alternating sum
is the stated determinant.

For the tableau formula, use \(r\ge1\) and take a finite rectangular
piece of the integer grid containing all starts \(A_i=(-i,1)\) and
ends \(B_j=(\lambda_j-j,r)\). Use east edges of weight \(x_v\) at
height \(v\) and north edges of weight \(1\). The path sum from \(A_i\)
to \(B_j\) is \(h_{\lambda_j-j+i}\): the heights of the east steps are
weakly increasing and determine a unique path. In zero horizontal
distance there is a single north path; negative distance gives none.
The matrix is the transpose of the Jacobi--Trudi matrix. The weighted
path determinant of \cite[Chapter 3]{finite-combinatorics} writes its
determinant as the signed sum over vertex-disjoint families. The grid
is acyclic, and the starts and the ends are each distinct. Both
boundary lines order endpoints from right to left by their index.
An inversion of that order forces intersection: at each height the
vertices of a path form an interval, and the north steps send its
right endpoint to the left endpoint of its next interval. Two disjoint
intervals preserve their order at every height, including the case
of one height. Hence only the identity matching survives.

For identity matching, the east-step heights in path \(i\) record
row \(i\) of a tableau. For adjacent paths let \(u_j,v_j\) be their
respective \(j\)-th heights. Where both rows have column \(j\), at
horizontal coordinate \(j-i-1\) the first path's vertical interval
is \([u_{j-1},u_j]\) and the second path's is \([v_j,v_{j+1}]\).
Use \(u_0=1\) and the final height \(r\) after the last step.
These are exactly the shared horizontal coordinates of adjacent
paths, since \(\lambda_i\ge\lambda_{i+1}\).
For \(j=1\) the first interval begins at \(1\); disjointness forces
\(u_1<v_1\). Inductively \(u_{j-1}<v_{j-1}\le v_j\), so at the next
coordinate disjointness again forces \(u_j<v_j\). The converse follows
from these same intervals, and their ordered separation prevents
intersection between nonadjacent paths as well: each intermediate
path is present at every shared horizontal coordinate of the two
outer paths. This gives the claimed weight-preserving bijection.
The argument works with any number of rows in the determinant once
it is at least the number of nonzero parts. A column taller than \(r\)
admits no strict filling. With no variables only \(h_0\) is nonzero;
the empty diagram gives \(1\), and for a nonempty shape the first row
of the determinant is zero. Compare \cite[Sections 4.4--4.6]{sagan}.
\end{proof}

Use the ring of symmetric functions in the following degreewise
sense. A degree \(n\) symmetric function is a compatible sequence of
homogeneous symmetric polynomials under setting the last variable to
zero. Its coefficients are determined by the monomial sums for
partitions of \(n\), which stabilize once there are at least \(n\)
variables. The tableau formula shows that \(s_\lambda\) is compatible
with this specialization. The products \(p_\nu\) are a basis over
\(\mathbb C\):
their monomial expansion first assigns the parts of \(\nu\) to
distinct variables, with nonzero coefficient \(\prod_j m_j(\nu)!\);
all other terms merge parts and have fewer distinct variables.
This gives a triangular change of basis. Define the Hall Hermitian
inner product by \(\langle p_\nu,p_\mu\rangle
=\boldsymbol{1}_{\nu=\mu}z_\nu\), linear in its first variable,
with different degrees orthogonal.

\begin{theorem}[Frobenius characteristic]\label{res:C05}
For a class function \(f\) on \(S_n\), the correspondence
\[
 \operatorname{ch}(f)=\sum_{\nu\vdash n}f(\nu)p_\nu/z_\nu
\]
is a degreewise isometric isomorphism and sends \(\chi^\lambda\) to
\(s_\lambda\). For class functions on \(S_a\) and \(S_b\), the
characteristic of their external tensor product induced from
\(S_a\times S_b\) to \(S_{a+b}\) is the product of their
characteristics.
\end{theorem}

\begin{proof}
The class sizes in Result~\ref{res:A01} and the defining norms of the
power sums give the isometry and bijection. For a partition \(\nu\),
the polynomial \(\Delta p_\nu\) is alternating. Its coefficients on
strictly decreasing exponent lists determine it, since lists with a
repeated exponent have zero coefficient. By Result~\ref{res:C03},
expansion in the \(a_{\lambda+\delta}\) and division by \(\Delta\)
give \(p_\nu=\sum_{\lambda}\chi^\lambda(\nu)s_\lambda\) in enough
variables. Orthogonality, as used in Result~\ref{res:C03}, inverts
this matrix and gives the claimed image of \(\chi^\lambda\).

For multiplication, a coset of \(S_a\times S_b\) is an ordered
decomposition into two label blocks, and the cosets fixed by a
permutation are assignments of whole cycles to the two blocks.
The trace on the summand for such a coset is \(f(\alpha)g(\beta)\),
where \(\alpha,\beta\) are the resulting cycle types. For fixed
types, the number of those assignments is
\(\prod_j\binom{m_j(\nu)}{m_j(\alpha)}=z_\nu/(z_\alpha z_\beta)\).
After division by \(z_\nu\), summing over the types is precisely the
product rule \(p_\alpha p_\beta=p_{\alpha\cup\beta}\). The trace
argument proves this for characters, hence by linearity for all
class functions. See \cite[Sections 4.6--4.7]{sagan}.
\end{proof}

\begin{corollary}[Young's rule]\label{res:C06}
The multiplicity of \(S^\lambda\) in the Young permutation module of
content \(\mu\) is \([x^\mu]s_\lambda(x)\). This is the number of
semistandard tableaux of shape \(\lambda\) and content \(\mu\).
\end{corollary}

\begin{proof}
Take the inner product of \(\xi^\mu\) with \(\chi^\lambda\).
By Result~\ref{res:C01} its value is the coefficient of \(x^\mu\) in
\(\sum_\nu\chi^\lambda(\nu)p_\nu(x)/z_\nu\), since the irreducible
character is real by Result~\ref{res:C03}. This polynomial is
\(s_\lambda\) by Result~\ref{res:C05}. Apply the tableau formula in
Result~\ref{res:C04}. See \cite[Theorem 2.11.2]{sagan}.
\end{proof}

\Needspace{7\baselineskip}
\section{Pieri and border-strip rules}

\begin{theorem}[Pieri rules]\label{res:C07}
The product \(s_\lambda h_k\) is the sum of \(s_\mu\) over
\(\mu\supseteq\lambda\) for which \(\mu/\lambda\) has \(k\) cells
with at most one in each column. For \(s_\lambda e_k\), the
corresponding condition is at most one in each row. Each summand
has coefficient \(1\). By the characteristic map these are the
decompositions obtained by inducing the external tensor product of
\(S^\lambda\) with the trivial or sign representation of \(S_k\),
respectively.
\end{theorem}

\begin{proof}
Work in enough variables for the degree in question and put
\(b=\lambda+\delta\). The horizontal-strip condition says
\(\mu_j\ge\lambda_j\) for every \(j\) and
\(\mu_j\le\lambda_{j-1}\) for \(j\ge2\): if that inequality fails,
column \(\lambda_{j-1}+1\) gains cells in rows \(j-1\) and \(j\).
Conversely, if column \(q\) gains a cell in row \(j\), then
\(q\le\mu_j\le\lambda_{j-1}\), so every earlier cell of that column
was already present. In strict
exponents \(c=\mu+\delta\), the independent ranges are \(c_1\ge b_1\)
and \(b_j\le c_j<b_{j-1}\) for \(j\ge2\). Therefore
\(\sum_c t^{\sum_j(c_j-b_j)}a_c(x)\) is the determinant obtained by
summing these geometric series in each column separately. Extract
\((1-x_i t)^{-1}\) from row \(i\). The new first column is \(x_i^{b_1}\);
column \(j\ge2\) is \(x_i^{b_j}-t^{b_{j-1}-b_j}x_i^{b_{j-1}}\).
The column-change matrix from \([x_i^{b_j}]\) is triangular with
diagonal \(1\). The determinant is thus
\(a_b(x)\prod_i(1-x_it)^{-1}\). Extracting \(t^k\) and dividing by
\(\Delta\) gives the first formula.

For the vertical-strip formula, use the alternating sum of
monomials defining \(a_b\). Multiplication by the symmetric
polynomial \(e_k\) commutes with alternation. Before alternating,
each term increments \(k\) different exponents by \(1\). A repeated
exponent gives zero, and the remaining lists stay strictly decreasing
because adjacent exponents originally differed by at least \(1\).
They give exactly the partitions with one added cell in each of
\(k\) distinct rows. Finally a single row has Schur function \(h_k\)
and a single column has \(e_k\) by Result~\ref{res:C04}. The last
assertion follows from Results~\ref{res:A07} and \ref{res:C05}.
The case \(k=1\) agrees with Result~\ref{res:B05}.
Compare \cite[Section 4.9]{sagan}.
\end{proof}

\begin{theorem}[Murnaghan--Nakayama rule]\label{res:C08}
For \(m\ge1\) and \(\lambda\vdash n\) with \(m\le n\), let
\(\tau\vdash n-m\). Then
\[
 \chi^\lambda((m)\cup\tau)
   =\sum_{\substack{\mu\subseteq\lambda\\
             \lambda/\mu\ \text{a border strip of size }m}}
        (-1)^{\operatorname{ht}(\lambda/\mu)}\chi^\mu(\tau).
\]
A border strip is connected by edges and contains no \(2\)-by-\(2\)
square. Its height is its number of occupied rows minus one.
The empty character takes value \(1\) on the empty cycle type.
\end{theorem}

\begin{proof}
For a partition \(\mu\), set \(b=\mu+\delta\). Multiplying the
alternating monomial sum for \(a_b\) by \(p_m\) adds \(m\) to one
exponent \(b_j\) before alternation. Discard coincident exponents.
Move the raised exponent to its place \(i\le j\) in the strictly
decreasing list \(c\); each interchange contributes a sign, giving
overall \((-1)^{j-i}\). The new partition \(\lambda=c-\delta\)
agrees with \(\mu\) outside rows \(i,\ldots,j\), increases in row
\(i\), and has
\(\lambda_k=\mu_{k-1}+1\) for \(i<k\le j\). All these rows gain
cells, and consecutive added row intervals meet in exactly one
column. Thus the skew shape is a border strip of height \(j-i\).
Conversely, for a border strip with occupied rows \(i,\ldots,j\),
connectedness gives an overlap at every adjacent pair of rows.
Because there is no square that overlap has length one.
The added intervals start at \(\mu_k+1\), so this forces
\(\lambda_k=\mu_{k-1}+1\) for \(i<k\le j\).
The total of \(m\) added cells determines \(\lambda_i\), recovering
exactly the exponent move above. Hence
\(p_m s_\mu=\sum_\lambda(-1)^{\operatorname{ht}(\lambda/\mu)}s_\lambda\).

The coefficient of \(s_\lambda\) in \(p_{(m)\cup\tau}=p_m p_\tau\)
is \(\chi^\lambda((m)\cup\tau)\) by the expansion in
Result~\ref{res:C05}. Expand \(p_\tau\) first and apply the
edge-strip multiplication just proved. This gives the rule,
including the initial value. See \cite[Section 4.10]{sagan}.
\end{proof}

\paragraph{Example: the character table of \(S_4\).}
Class sizes follow from Result~\ref{res:A01}. In the standard
representation the trace is the number of fixed points minus one
by Result~\ref{res:B09}. Its sign twist gives the row \((2,1,1)\).
For \((2,2)\), a \(2\)-strip leaves either a horizontal pair with
positive sign or a vertical pair with negative sign. A \(3\)-strip
leaves one cell with negative sign, and the full square is no
\(4\)-strip. The value at the identity is \(2\) by
Result~\ref{res:B08}. These computations give
\[
\begin{array}{c|rrrrr}
 \nu & 1^4 & (2,1,1) & (2,2) & (3,1) & (4)\\ \hline
 \text{class size} & 1 & 6 & 3 & 8 & 6\\ \hline
 (4)     &1& 1& 1& 1& 1\\
 (3,1)   &3& 1&-1& 0&-1\\
 (2,2)   &2& 0& 2&-1& 0\\
 (2,1,1) &3&-1&-1& 0& 1\\
 (1^4)   &1&-1& 1& 1&-1
\end{array}
\]
With the indicated class-size weights the row Gram matrix is
\(24I\). The squared column norms are \(24,4,8,3,4\), the five
centralizer orders, and distinct columns are orthogonal.

\Needspace{26\baselineskip}
\section{The dual Jacobi--Trudi formula}

\begin{corollary}[Dual Jacobi--Trudi]\label{res:C09}
Let \(\lambda^{\mathsf t}\) be the transpose of the partition
\(\lambda\), and pad it with zeros to length \(q\ge \lambda_1\)
with \(\lambda_1=0\) for the empty partition.
Then in any finite alphabet \(x\),
\[
 s_\lambda(x)=
 \det\bigl[e_{\lambda^{\mathsf t}_i-i+j}(x)\bigr]_{i,j=1}^{q}.
\]
Here \(e_0=1\) and \(e_k=0\) for \(k<0\). For the empty shape one
may take \(q=0\) and use the empty determinant \(1\).
\end{corollary}

\begin{proof}
On the stable symmetric functions used in \ref{res:C05}, define
\(\omega(p_j)=(-1)^{j-1}p_j\), extending multiplicatively and
linearly. The products \(p_\nu\) are a basis, so this rule defines
an algebra involution. Both generating series are identities
in stable symmetric functions:
\begin{align*}
 \sum_{k\ge0}h_k t^k &= \exp\!\left(\sum_{j\ge1}p_jt^j/j\right),\\
 \sum_{k\ge0}e_k t^k &= \exp\!\left(\sum_{j\ge1}(-1)^{j-1}p_jt^j/j\right).
\end{align*}
In a finite alphabet these follow by taking formal logarithms of
\(\prod_i(1-x_it)^{-1}\) and \(\prod_i(1+x_it)\), respectively.
Thus \(\omega(h_k)=e_k\), also with the zero convention in
negative degree.

On a permutation of cycle type \(\nu\), the sign is
\(\prod_{j\in\nu}(-1)^{j-1}\), counting parts with multiplicity.
The characteristic map of \ref{res:C05} therefore intertwines
tensoring with the sign representation and applying \(\omega\).
By \ref{res:A09} it follows that
\(\omega(s_\mu)=s_{\mu^{\mathsf t}}\). Apply \(\omega\) to the
ordinary Jacobi--Trudi formula of \ref{res:C04} for
\(\mu=\lambda^{\mathsf t}\). Use \(q\) rows in the determinant
before applying \(\omega\), and then specialize the stable
identity to the chosen finite alphabet. This gives the formula
for all alphabet sizes. See \cite[Theorem 4.5.1]{sagan}.
\end{proof}

\chapter{Schur--Weyl Duality}

Let \(V\) be a complex vector space of dimension \(d\), and put
\(X=V^{\otimes n}\). The two actions on \(X\) exchange the tensor
positions and change the vectors in all positions simultaneously.
Their images in \(\operatorname{End}(X)\) are each other's centralizers.
We first prove the two centralizer identities, then determine the
summands and their characters.

\section{Commuting actions and centralizers}

\begin{proposition}[The two tensor actions]
\label{res:D01}
The symmetric group \(S_n\) and the group \(\operatorname{GL}(V)\)
act linearly on \(X\) by
\begin{align*}
 P_\sigma(v_1\otimes\cdots\otimes v_n)
   &=v_{\sigma^{-1}(1)}\otimes\cdots\otimes v_{\sigma^{-1}(n)},\\
 G_g(v_1\otimes\cdots\otimes v_n)
   &=gv_1\otimes\cdots\otimes gv_n .
\end{align*}
They commute. Thus the images
\[
 A=\operatorname{im}\bigl(\mathbb C[S_n]\longrightarrow
                         \operatorname{End}(X)\bigr),\qquad
 B=\operatorname{span}_{\mathbb C}\{G_g:g\in\operatorname{GL}(V)\}
\]
are commuting subalgebras of \(\operatorname{End}(X)\).
\end{proposition}

\begin{proof}
The displayed expressions are multilinear in the input vectors.
Under successive permutations \(\tau\) and \(\sigma\), the input
at output position \(i\) comes from \(\tau^{-1}\sigma^{-1}(i)\).
This proves \(P_\sigma P_\tau=P_{\sigma\tau}\), and the identity
permutation acts as the identity. The group law
\(G_gG_h=G_{gh}\) holds in every position. The two orders of action
give \(gv_{\sigma^{-1}(i)}\) at position \(i\), so they agree on
the spanning pure tensors. Extending the permutation action
linearly gives the algebra map defining \(A\). Products of the
generators of \(B\) are again generators, and \(G_1\) is its identity.

For \(n=0\) we have \(X=\mathbb C\) with both actions trivial.
For \(d=0<n\) we have \(X=0\), whose endomorphism algebra is the
zero algebra. We use the unique element of that algebra as its
identity. These conventions apply throughout this chapter.
The actions agree with those of \cite[\S\S4.18--4.19]{etingof}.
\end{proof}

\begin{lemma}[Endomorphisms as tensors]
\label{res:D02}
The map
\[
 \Phi\colon\operatorname{End}(V)^{\otimes n}\longrightarrow
                 \operatorname{End}(X),\qquad
 T_1\otimes\cdots\otimes T_n\longmapsto
 \bigl(v_1\otimes\cdots\otimes v_n\mapsto T_1v_1\otimes\cdots\otimes T_nv_n\bigr)
\]
is an algebra isomorphism. It identifies conjugation by \(P_\sigma\)
with permutation of the tensor positions in the source.
Consequently \(\operatorname{End}_{A}(X)\) corresponds to the
\(S_n\)-invariant tensors.
\end{lemma}

\begin{proof}
For \(d,n\geq1\), choose a basis \(v_1,\ldots,v_d\) of \(V\)
and let \(E_{ab}\) send \(v_b\) to \(v_a\) and all other basis
vectors to zero. With \(\mathbf a,\mathbf b\in\{1,\ldots,d\}^n\),
\(E_{a_1b_1}\otimes\cdots\otimes E_{a_nb_n}\) maps to the matrix
unit in the tensor basis taking the word \(\mathbf b\) to
\(\mathbf a\). Thus \(\Phi\) takes a basis to a basis.
Composition of pure tensor operators is multiplication position
by position, and the tensor of the identities is the identity.
Evaluation of \(P_\sigma\Phi(T_1\otimes\cdots\otimes T_n)P_{\sigma^{-1}}\)
on a pure tensor gives \(T_{\sigma^{-1}(i)}\) at position \(i\).
An operator commutes with \(A\) exactly when it commutes with
every \(P_\sigma\), which proves the invariance criterion.
For \(n=0\) the map is the scalar identification, and for \(d=0<n\)
both vector spaces are zero. This is the comparison used in
\cite[Theorem 4.55]{etingof}.
\end{proof}

\begin{theorem}[The first centralizer identity]
\label{res:D03}
The centralizer of \(A\) in \(\operatorname{End}(X)\) is \(B\).
\end{theorem}

\begin{proof}
For \(n=0\) the two algebras and the ambient endomorphism algebra
are \(\mathbb C\), and for \(d=0<n\) all are zero. Suppose \(d,n\geq1\)
and set \(W=\operatorname{End}(V)\). In \(W^{\otimes n}\) every
invariant tensor is its own average over \(S_n\). Since pure
tensors span, the invariant tensors are spanned by sums of
permutations of pure tensors. For \(T_1,\ldots,T_n\in W\),
multilinearity gives
\begin{equation*}
 \sum_{J\subseteq\{1,\ldots,n\}} (-1)^{n-|J|}
      \biggl(\sum_{j\in J}T_j\biggr)^{\otimes n}
  =\sum_{\sigma\in S_n}
         T_{\sigma(1)}\otimes\cdots\otimes T_{\sigma(n)} .
\end{equation*}
Indeed a word in indices using exactly the set \(I\) occurs on the
left with coefficient \(\sum_{J\supseteq I}(-1)^{n-|J|}\).
This coefficient is zero for a proper subset \(I\) and is one for
the full set. A word of length \(n\) using the full set uses each
index once. The invariant tensors are therefore spanned by the
powers \(T^{\otimes n}\), since these powers are invariant.

For any \(T\in W\), the determinant of \(T+tI\) is a monic
polynomial of degree \(d\). There exist \(n+1\) distinct
\(t_0,\ldots,t_n\) for which \(T+t_jI\) is invertible.
The vector-valued polynomial \(t\mapsto(T+tI)^{\otimes n}\)
has degree at most \(n\). The scalar Lagrange interpolation
formula, applied in each coordinate, expresses its value at
zero as a linear combination of its values at the chosen
points. Under \(\Phi\) those values are generators of \(B\).
Lemma~\ref{res:D02} now gives the required centralizer.
Compare \cite[Lemma 4.56 and Proposition 4.58]{etingof}.
\end{proof}

\begin{theorem}[The second centralizer and the tensor decomposition]
\label{res:D04}
Write
\[
 L_\lambda(V)=\operatorname{Hom}_{S_n}(S^\lambda,X),
 \qquad (g\cdot f)(u)=G_g f(u)
 \quad (g\in\operatorname{GL}(V)).
\]
Evaluation is an isomorphism of \(S_n\times\operatorname{GL}(V)\)
representations
\[
 \bigoplus_{\lambda\vdash n} S^\lambda\otimes L_\lambda(V)
        \longrightarrow X .
\]
Every nonzero \(L_\lambda(V)\) is irreducible for
\(\operatorname{GL}(V)\), and these nonzero modules are mutually
nonisomorphic. The centralizer of \(B\) is \(A\).
\end{theorem}

\begin{proof}
By Maschke and Schur \cite[Chapter 8]{algebra-core} and the
classification in \ref{res:A07}, \(X\) is a direct sum of copies
of the \(S^\lambda\). On a copy of \(S^\lambda\), evaluation
from \(S^\lambda\otimes\operatorname{Hom}_{S_n}(S^\lambda,S^\lambda)\)
is scalar evaluation and hence an isomorphism. The Hom spaces
between different simple modules vanish. Applying these facts to
each copy proves the evaluation isomorphism. Postcomposition
defines the asserted action on \(L_\lambda(V)\) because \(G_g\)
commutes with \(S_n\). Evaluation is equivariant for both actions.

Put \(R=\mathbb C[S_n]\). We first identify its full image on
the simple modules. The action map
\[
 R\longrightarrow \prod_{\lambda\vdash n}\operatorname{End}(S^\lambda)
\]
is injective because any element that annihilates every simple
module annihilates the semisimple left regular module and hence
annihilates \(1\) in that module. The multiplicity of \(S^\lambda\)
in the regular module is \(\dim S^\lambda\). Indeed evaluation
at \(1\) identifies \(\operatorname{Hom}_R(R,S^\lambda)\) with
\(S^\lambda\), and computing the same Hom space with the
regular module decomposed into simple summands gives the
multiplicity by Schur. It follows that
\(\dim R=\sum_{\lambda}(\dim S^\lambda)^2\), so the injective
action map is an isomorphism.

Across the summands \(S^\lambda\otimes L_\lambda(V)\) with
\(L_\lambda(V)\ne0\), this identifies \(A\) with the product of the
independent full matrix blocks \(\operatorname{End}(S^\lambda)\).
Each block acts on the first factor and by the identity on the
second. In particular \(A\) contains the projectors to these
summands. An operator commuting with \(A\) preserves every
summand because it commutes with those projectors. In a basis
\((u_i)\) of \(S^\lambda\), commuting with each diagonal matrix
unit makes the operator preserve every \(u_i\otimes L_\lambda(V)\).
Commuting with each off-diagonal matrix unit identifies the
operators on these spaces with one common
\(T\in\operatorname{End}(L_\lambda(V))\). Conversely
\(I\otimes T\) commutes with the entire first block. By
Theorem~\ref{res:D03}, \(B\) is the product of these independent
second-factor blocks.

An operator commuting with \(B\) preserves the same summands
because \(B\) also contains their projectors. Applying the matrix
unit argument to the second factor shows that the new centralizer
is \(A\). Since the linear span of the group on \(L_\lambda(V)\)
is the full matrix algebra, a nonzero invariant subspace contains
the image of each of its vectors under all linear maps, so it is the
whole space. A group intertwiner between two such spaces intertwines
the whole span \(B\), whose projector onto one summand is the
identity there and zero on the other. It must be zero.
If \(X=0\), the evaluation and centralizer identities are between
zero spaces or algebras. The block argument gives the
double centralizer conclusion of
\cite[Theorems 4.54 and 4.57 and Corollary 4.59]{etingof}.
\end{proof}

\section{The length bound and the group algebra kernel}

\begin{theorem}[The exact length bound]
\label{res:D05}
For \(\lambda\vdash n\), \(L_\lambda(V)\) is nonzero if and only
if the number of nonzero parts \(\ell(\lambda)\) is at most \(d\).
\end{theorem}

\begin{proof}
For a set \(C\) of positions let \(b_C\in\mathbb C[S_n]\) be the
sum of the permutations of \(C\), weighted by their signs and
extended by the identity off \(C\). If \(|C|>d\), each basis
tensor in \(X\) has a repeated basis vector among those positions.
The transposition of these two positions fixes that tensor.
Pairing permutations of \(C\) that differ on the right by this
transposition shows that \(b_C\) sends the tensor to zero.
The same conclusion holds if there are no basis tensors because
\(X=0\).

If \(\ell(\lambda)>d\), take \(C\) to be the first column of a
tableau \(t\) of shape \(\lambda\). The polytabloid \(e_t\) in
\ref{res:A04} transforms
by the sign under each permutation of \(C\), so
\(b_C e_t=|C|!\,e_t\). An equivariant map from \(S^\lambda\) to
\(X\) sends this multiple to zero.
The scalar is nonzero, and \(e_t\) generates \(S^\lambda\), so
the map itself is zero.

Suppose \(\ell(\lambda)\leq d\). For a tabloid \(\{t\}\), take the
tensor with the \(i\)-th basis vector of \(V\) at every position whose
label belongs to row \(i\) of \(t\). This construction is unchanged
by reordering entries within a row. Distinct tabloids give
distinct tensor basis words, since the word recovers each row's
label set. Formula \ref{res:D01} for the left action shows that
relabeling the tabloid permutes tensor positions equivariantly.
These words span the subspace of \(X\) of content \(\lambda\)
in the chosen basis, and the tabloid map identifies that subspace
with \(M^\lambda\). It therefore restricts to an injection
of \(S^\lambda\) by \ref{res:A04}. The empty shape for \(n=0\)
gives the empty word, a basis of \(X=\mathbb C\).
The length criterion also appears in
\cite[Theorem 4.63]{etingof}, where it follows from the character.
\end{proof}

\begin{corollary}[The alternating generator of the kernel]
\label{res:D06}
In the full matrix block decomposition of \(\mathbb C[S_n]\),
the kernel of its action on \(X\) is the sum of the blocks with
\(\ell(\lambda)>d\). For \(n>d\) this is the two-sided ideal
generated by the alternating sum \(b_{\{1,\ldots,d+1\}}\).
For \(n\leq d\) the action of the group algebra is injective.
\end{corollary}

\begin{proof}
Theorem~\ref{res:D04} identifies the action with the product of
the blocks with \(L_\lambda(V)\ne0\). Theorem~\ref{res:D05}
identifies those omitted. For \(n>d\), the pairing argument
of Theorem~\ref{res:D05} shows that the indicated alternating
sum acts as zero on \(X\). In each omitted block choose a
tableau \(t\) with \(1,\ldots,d+1\) in its first column.
The alternating sum sends \(e_t\) to \((d+1)!\,e_t\), so its
projection to that block is nonzero. Central block projectors
separate the block components of the two-sided ideal.
A nonzero matrix \(T\) generates the entire two-sided ideal in a
full matrix algebra, since \(E_{ia}TE_{bj}=T_{ab}E_{ij}\) for
a nonzero entry \(T_{ab}\). The ideal is therefore exactly
the sum of all omitted blocks. For \(d=0<n\) the alternating
generator on one position is the unit, as required.
For \(n\leq d\) all partitions of \(n\) have length at most \(d\).
\end{proof}

\section{Characters and dimensions}

\begin{proposition}[The trace along permutation cycles]
\label{res:D07}
For \(g\in\operatorname{GL}(V)\) and \(\sigma\in S_n\),
\[
 \operatorname{tr}_X(P_\sigma G_g)
       =\prod_{c\text{ a cycle of }\sigma}
                   \operatorname{tr}_V(g^{|c|}).
\]
If \(\chi^\lambda\) is the character of \(S^\lambda\), then
\begin{equation}
\label{eq:schur-weyl-trace-average}
 \operatorname{tr}_{L_\lambda(V)}(g)
 =\frac{1}{n!}\sum_{\sigma\in S_n}
     \chi^\lambda(\sigma^{-1})
     \prod_{c\text{ a cycle of }\sigma}\operatorname{tr}_V(g^{|c|}).
\end{equation}
\end{proposition}

\begin{proof}
For \(d,n\geq1\), let \(g_{ab}\) be the matrix entry in output
row \(a\) and input column \(b\). For a basis word with index
\(\mathbf i\), the corresponding diagonal matrix entry of
\(P_\sigma G_g\) is \(\prod_j g_{i_{\sigma(j)},i_j}\).
Summing over the indices divides into independent sums over the
cycles. On a cycle \((j_1,\ldots,j_k)\) the factors are
\(g_{i_{j_2},i_{j_1}}\cdots g_{i_{j_1},i_{j_k}}\), whose sum is
\(\operatorname{tr}_V(g^k)\) by the definition of matrix
multiplication and trace. For \(n=0\) the trace and the empty
product equal one. For \(d=0<n\) they both equal zero.

Using the evaluation decomposition of Theorem~\ref{res:D04},
the same trace is
\(\sum_{\mu\vdash n}\chi^\mu(\sigma)\operatorname{tr}_{L_\mu(V)}(g)\).
Finite-group character orthogonality
\cite[Chapter 8]{algebra-core} gives
\(\frac{1}{n!}\sum_\sigma\chi^\lambda(\sigma^{-1})\chi^\mu(\sigma)
=\delta_{\lambda\mu}\). To see the inverse in this formula,
every \(S_n\)-representation can be made unitary by averaging a
positive definite Hermitian form, so the character on the inverse
is the complex conjugate. Multiplying the decomposition by
\(\chi^\lambda(\sigma^{-1})\) and averaging yields
\eqref{eq:schur-weyl-trace-average}. See \cite[\S4.21]{etingof}.
\end{proof}

\begin{theorem}[Weyl character and dimension]
\label{res:D08}
For invertible \(g=\operatorname{diag}(x_1,\ldots,x_d)\),
\(\operatorname{tr}_{L_\lambda(V)}(g)=s_\lambda(x_1,\ldots,x_d)\),
giving the polynomial character on the diagonal.
If \(\ell(\lambda)\leq d\), pad \(\lambda\) with zeros to length
\(d\). Then
\[
 \dim L_\lambda(V)=
    \prod_{1\leq i<j\leq d}\frac{\lambda_i-\lambda_j+j-i}{j-i}.
\]
For \(\ell(\lambda)>d\) the dimension is zero.
\end{theorem}

\begin{proof}
On the diagonal, \(\operatorname{tr}_V(g^k)=p_k(x)\).
Each permutation is conjugate to its inverse because the two
have the same cycle lengths (\ref{res:A01}). The cycle
class \(\nu\) has \(n!/z_\nu\) elements. Grouping terms in
\eqref{eq:schur-weyl-trace-average} yields
\[
 \operatorname{tr}_{L_\lambda(V)}(g)
    =\sum_{\nu\vdash n}\frac{\chi^\lambda(\nu)}{z_\nu}p_\nu(x)
    =s_\lambda(x)
\]
by the Frobenius characteristic identity \ref{res:C05}. This
also gives the polynomial extension of the diagonal character.
For too many rows the semistandard model \ref{res:C04}
has an empty set of fillings with entries from \(1,\ldots,d\),
in agreement with Theorem~\ref{res:D05}.

For \(d\geq1\) and \(\ell(\lambda)\leq d\), use the
alternant ratio in \ref{res:C04}. Write \(b_j=\lambda_j+d-j\)
and \(a_j=d-j\), both strictly decreasing sequences.
Set \(x_i=t^{d-i}\) for real \(t>1\). The numerator matrix
has entries \((t^{b_j})^{d-i}\), a Vandermonde matrix with rows
indexed by decreasing powers. The denominator is the same
matrix with \(b_j\) replaced by \(a_j\). Thus the ratio is
\(\prod_{i<j}(t^{b_i}-t^{b_j})/(t^{a_i}-t^{a_j})\).
Each factor tends to \((b_i-b_j)/(a_i-a_j)\) as \(t\to1\),
by dividing both differences by \(t-1\). The Schur polynomial
is continuous at \((1,\ldots,1)\), where the trace is the
dimension. This gives the formula. For \(d=0\), the only
permitted shape is empty and its space is \(\mathbb C\).
For the empty shape every factor is one, and for \(d=0,1\)
the product has no factors. These are the formulas of
\cite[Theorem 4.63]{etingof}.
\end{proof}

\begin{example}[Symmetric and exterior powers]
\label{res:D09}
For \(n\geq1\) the row shape and column shape give, respectively,
\[
 L_{(n)}(V)\cong\operatorname{Sym}^n V,\qquad
 L_{(1^n)}(V)\cong\bigwedge^n V.
\]
The zero-degree conventions give \(\mathbb C\) in both cases.
If \(d=2\) and \(n=3\), the tensor decomposition has dimensions
\(8=4+2\cdot2\), with no summand from the shape \((1,1,1)\).
\end{example}

\begin{proof}
By \ref{res:A07} the two simple \(S_n\)-modules are the
trivial and sign representations. Evaluating their intertwining
maps at \(1\) identifies their multiplicity spaces with the
invariant and alternating subspaces, respectively. For
\(\epsilon=1\) or \(\epsilon=\operatorname{sgn}\), set
\[
 Q_\epsilon=\frac{1}{n!}\sum_{\sigma\in S_n}\epsilon(\sigma)P_\sigma .
\]
Reindexing the sum by left or right multiplication gives
\(P_\tau Q_\epsilon=Q_\epsilon P_\tau=\epsilon(\tau)Q_\epsilon\).
Also \(Q_\epsilon\) restricts to the identity on the vectors
transforming by \(\epsilon\). Its kernel is the subspace
spanned by \(v-\epsilon(\sigma)P_\sigma v\), since it kills
these differences and \(v-Q_\epsilon v\) is their average.
Thus it identifies the isotypic subspace with the corresponding
quotient. For \(\epsilon=1\) that quotient is the symmetric
power by the multilinear definition. For the sign it is the exterior
power. An alternating multilinear map kills the signed
permutation differences, by exchanging adjacent positions.
Conversely in the quotient a transposition gives the sign rule and
a tensor with two equal positions is its own negative, hence
zero. These properties give precisely the universal
alternating multilinear map. Both identifications commute
with the diagonal \(\operatorname{GL}(V)\) action.

When \(d=2,n=3\), Theorem~\ref{res:D05} leaves only the
shapes \((3)\) and \((2,1)\). Their Specht dimensions are
\(1,2\) (the standard basis or the hook-length formula
\ref{res:B08}). Theorem~\ref{res:D08} gives multiplicity-space
dimensions \(4,2\). This proves the dimension check and
recovers the symmetric and alternating cases of
\cite[Example 4.60]{etingof}.
\end{proof}

\section{Lie algebra and polynomial representations}

\begin{corollary}[The Lie algebra action]
\label{res:D10}
For \(T\in\mathfrak{gl}(V)=\operatorname{End}(V)\) with the
commutator bracket, let
\[
 \mathcal D(T)=\sum_{i=1}^n
       I^{\otimes(i-1)}\otimes T\otimes I^{\otimes(n-i)}
       \quad\text{acting on }X.
\]
The induced action of its universal enveloping algebra has
image \(B\). In particular its action has the same centralizer
\(A\), and the nonzero multiplicity spaces in
Theorem~\ref{res:D04} are irreducible and mutually
nonisomorphic as \(\mathfrak{gl}(V)\)-modules.
\end{corollary}

\begin{proof}
On distinct positions the summand operators commute.
On a common position their commutator acts by \([T,U]\).
This gives \([\mathcal D(T),\mathcal D(U)]=\mathcal D([T,U])\).
Extending the linear map to the tensor algebra and imposing
the defining commutator relations yields a unital algebra map
from the universal enveloping algebra. Write \(B'\) for its image.

For \(d,n\geq1\), \(\mathcal D(T)\) is the coefficient of \(t\)
in \((I+tT)^{\otimes n}\). Choose \(n+1\) distinct \(t\)'s
where \(I+tT\) is invertible. Interpolation at these points
expresses that coefficient as a linear combination of operators
in \(B\). This gives \(B'\subseteq B\).

Every invertible complex matrix is a product of exponentials.
Gaussian elimination expresses an invertible matrix as a
product of coordinate shears, invertible diagonal matrices
and coordinate swaps. A shear \(I+aE_{ij}\) for \(i\ne j\)
is \(\exp(aE_{ij})\), and an invertible diagonal matrix is
the exponential of the diagonal matrix of chosen scalar
logarithms. On two coordinates a swap matrix is a product of
three shears and a diagonal matrix, as seen from
\[
 \begin{pmatrix}1&1\\0&1\end{pmatrix}
 \begin{pmatrix}1&0\\-1&1\end{pmatrix}
 \begin{pmatrix}1&1\\0&1\end{pmatrix}
 \begin{pmatrix}-1&0\\0&1\end{pmatrix}
     =\begin{pmatrix}0&1\\1&0\end{pmatrix}.
\]
The summands defining \(\mathcal D(T)\) commute, so multiplying
their absolutely convergent exponential series gives
\(\exp(\mathcal D(T))=(\exp T)^{\otimes n}\). The exponential
belongs to \(B'\) because its finite sums belong to the unital
subalgebra \(B'\), and every finite-dimensional linear subspace
is closed. Hence \(G_g\in B'\) for every \(g\in\operatorname{GL}(V)\).
This proves \(B=B'\). For \(n=0\), the derived operators vanish
and the unital image is the scalar algebra, and for \(d=0<n\)
both images are zero. The centralizer and simplicity assertions
now follow from the block argument of Theorem~\ref{res:D04},
also proving the Lie algebra version in
\cite[Theorems 4.55 and 4.57]{etingof}.
\end{proof}

\begin{theorem}[Polynomial representations of the general linear group]
\label{res:D11}
A finite-dimensional representation of \(\operatorname{GL}(V)\)
is called polynomial if its matrix entries in a basis are
polynomials in the matrix entries of \(g\). Every such
representation splits into homogeneous degrees. At any fixed
degree \(n\), the representations are completely reducible with
simple modules exactly \(L_\lambda(V)\) for \(\lambda\vdash n\)
with \(\ell(\lambda)\leq d\).
\end{theorem}

\begin{proof}
First suppose \(d\geq1\) and let \(W\) carry a polynomial
representation \(\rho\). The scalar subgroup gives a finite
expansion \(\rho(tI)=\sum_{k=0}^{N}t^k P_k\) for \(t\ne0\).
Comparing coefficients in the polynomial identity
\(\rho(tI)\rho(sI)=\rho(tsI)\) for \(t,s\ne0\) gives
\(P_iP_j=0\) for distinct \(i,j\) and \(P_i^2=P_i\).
Their sum is \(I_W\) by setting \(t=1\), so their images form a
direct sum. Each \(P_k\) commutes with \(\rho(g)\), as follows
by comparing coefficients in the commutation identity for
scalar matrices. On \(W_k=\operatorname{im}P_k\) we have
\(\rho(tg)|_{W_k}=t^k\rho(g)|_{W_k}\). The matrix entries of the
restricted action are polynomials obtained by a fixed linear
inclusion and projection. Decomposing them by degree and
comparing coefficients of \(t\) shows that only the terms of
degree \(k\) contribute on \(\operatorname{GL}(V)\). This gives
the homogeneous degree decomposition.

Suppose that \(\rho\) is homogeneous of degree \(n\).
The matrix entries of \(G_g\) on the tensor basis are the
products \(\prod_j g_{a_jb_j}\). Every monomial of total degree
\(n\) occurs in this way, including the constant at \(n=0\).
If \(\sum_j c_jG_{g_j}=0\) in \(B\), each of these monomials
has weighted sum zero over the \(g_j\). Every entry of \(\rho\)
is a linear combination of the monomials, so
\(\sum_j c_j\rho(g_j)=0\). Thus \(G_g\mapsto\rho(g)\) extends
to a well-defined linear map from \(B\) to
\(\operatorname{End}(W)\). It respects multiplication because
the group actions do, and sends \(G_1\) to the identity.
This makes \(W\) a \(B\)-module. Conversely every
\(L_\lambda(V)\ne0\) has a homogeneous polynomial action of
degree \(n\): use the evaluation isomorphism to embed it in
\(X\) by tensoring with a fixed nonzero vector of \(S^\lambda\).
Composition with a fixed linear left inverse reads off
its matrix entries from those of \(G_g\).

By Theorems~\ref{res:D04} and \ref{res:D05}, \(B\) is the
product of the full matrix algebras on these spaces. Its
central projectors split every module into its parts for the
single blocks. To see the structure for one block
\(\operatorname{Mat}_{m}(\mathbb C)\), take a basis of the image
of \(E_{11}\) in a module. For each basis vector \(w\), the
vectors \(E_{i1}w\) form a copy of the standard module by the
matrix unit multiplication rule. These vectors for all \(w\)
are independent, as applying \(E_{1i}\) recovers the
coefficients of the \(i\)-th part. They span since
\(\sum_i E_{i1}E_{1i}=1\). Therefore the simple modules of \(B\)
are exactly the standard modules of the occurring blocks and
every module is a direct sum of them. Restricting to the group
gives the required result. For \(d=0\), the general linear
group is trivial and the only polynomials in the empty set of
matrix entries are constants. All nonzero representations
have degree zero, where the unique simple module is
\(L_{\varnothing}(0)=\mathbb C\).
\end{proof}

This is the matrix-entry polynomial case of
\cite[Theorem 4.66(i)]{etingof}. The later algebraic-group topic
treats rational representations (allowing \(\det^{-1}\)),
determinant twists, and algebraic Peter--Weyl
\cite[Theorem 4.66(ii)]{etingof}.


\backmatter
\begin{thebibliography}{9}
\addcontentsline{toc}{chapter}{Bibliography}

\bibitem{sagan}
Bruce E. Sagan.
\emph{The Symmetric Group: Representations, Combinatorial Algorithms,
and Symmetric Functions}, second edition.
Graduate Texts in Mathematics 203, Springer, 2001.
\href{https://doi.org/10.1007/978-1-4757-6804-6}{DOI}.

\bibitem{etingof}
Pavel Etingof et al.
\emph{Introduction to Representation Theory}.
MIT 18.712 lecture notes, Fall 2010, archived OCW version.
\href{https://ocw.mit.edu/courses/18-712-introduction-to-representation-theory-fall-2010/resources/mit18_712f10_replect/}{MIT OpenCourseWare}.

\bibitem{algebra-core}
MathGenome.
\emph{Algebra Core} (P1-03), accepted notes and formalization.
Chapters 8--9 treat ordinary finite-group representations,
characters, induction, and Mackey theory. MathGenome v0.21.0.

\bibitem{finite-combinatorics}
MathGenome.
\emph{Finite Combinatorics} (P2-04), accepted notes and formalization.
Chapter 3 includes the weighted Lindstr\"om--Gessel--Viennot path
determinant. MathGenome v0.21.0.
\end{thebibliography}
\end{document}
